% GENERATED FILE. Edit canonical lessons, then run npm run generate:books.
% canonical-source-hash: fnv1a64:33d726cf82b546c8
% canonical-lessons: ML-C358-pantupantu, ML-C358-akkalattu, ML-C358-anneram, ML-C358-tutarnnu, ML-C358-utanilam

\chapter{Once Upon a Time}
\label{ch:ml-once-upon-a-time}
\hlchaptermodality{358}

\begin{chapteropening}
\textbf{By the end of this chapter:} \emph{I can say once upon a time, in those days, at that moment, following that and throughout.}

Complete the last lesson of chapter 358: I can say once upon a time, in those days, at that moment, following that and throughout.
\end{chapteropening}

\section[\texorpdfstring{\ml{പണ്ടുപണ്ട്} (paṇṭupaṇṭŭ)}{പണ്ടുപണ്ട് (paṇṭupaṇṭŭ)}]{\ml{പണ്ടുപണ്ട്} (paṇṭupaṇṭŭ) — once upon a time}

\label{lesson:ML-C358-pantupantu}



\begin{quote}
Before the new one: say the Malayalam for ever since, then the Malayalam for all of a sudden.
\end{quote}

\subsection*{You'll want to know: \ml{പണ്ടുപണ്ട്}}
\textbf{\ml{പണ്ടുപണ്ട്}} — \emph{paṇṭupaṇṭŭ} — \textquotedblleft{}once upon a time\textquotedblright{}.

\begin{quote}
\ml{പണ്ടുപണ്ട്} \ml{ഒരു} \ml{ഗ്രാമത്തിൽ} \ml{ഒരു} \ml{പാവം} \ml{കർഷകൻ} \ml{ജീവിച്ചിരുന്നു}.
\end{quote}

\begin{quote}
\emph{paṇṭupaṇṭŭ oru grāmattil oru pāvaṁ karṣakan jīviccirunnu.}
\end{quote}

\begin{quote}
Once upon a time a poor farmer lived in a village.
\end{quote}

The doubled \ml{പണ്ട്}, long ago, opens a folk tale.

\begin{cousinweb}
\textbf{\ml{പണ്ടുപണ്ട്}} (\emph{paṇṭupaṇṭŭ}) doubles \ml{പണ്ട്} (\emph{paṇṭŭ}), \textquotedblleft{}long ago\textquotedblright{}, which DEDR 3999 sets with \ml{പഴയ}, \textquotedblleft{}old\textquotedblright{}: native Dravidian. Said twice, it pushes the story back into the time of tales.
\end{cousinweb}

\begin{grammarlens}[title={What you've built}]
One more word for telling what happened.
\end{grammarlens}

\subsection*{Your turn}
\begin{itemize}
\raggedright
  \item \emph{Say it:} \emph{paṇṭupaṇṭŭ}
  \item \emph{Say it:} \emph{paṇṭupaṇṭŭ}, once more
  \item \emph{Say it:} \emph{paṇṭupaṇṭŭ oru grāmattil oru pāvaṁ karṣakan jīviccirunnu.} — Once upon a time a poor farmer lived in a village.
  \item \emph{Recall:} say the Malayalam for ever since, then the Malayalam for all of a sudden, then say \emph{paṇṭupaṇṭŭ} again
\end{itemize}

\begin{tcolorbox}[breakable,colback=teal!4,colframe=teal!35!black,title={Before you move on}]
What does \ml{പണ്ടുപണ്ട്} mean? (\textquotedblleft{}Once upon a time\textquotedblright{}.) Say it once more.
\end{tcolorbox}
\section[\texorpdfstring{\ml{അക്കാലത്ത്} (akkālattŭ)}{അക്കാലത്ത് (akkālattŭ)}]{\ml{അക്കാലത്ത്} (akkālattŭ) — in those days}

\label{lesson:ML-C358-akkalattu}



\begin{quote}
Before the new one: say the Malayalam for all of a sudden, then the Malayalam for once upon a time.
\end{quote}

\subsection*{You'll want to know: \ml{അക്കാലത്ത്}}
\textbf{\ml{അക്കാലത്ത്}} — \emph{akkālattŭ} — \textquotedblleft{}in those days\textquotedblright{}.

\begin{quote}
\ml{അക്കാലത്ത്} \ml{ഗ്രാമത്തിൽ} \ml{ബസ്} \ml{ഉണ്ടായിരുന്നില്ല}.
\end{quote}

\begin{quote}
\emph{akkālattŭ grāmattil basŭ uṇṭāyirunnilla.}
\end{quote}

\begin{quote}
In those days there was no bus in the village.
\end{quote}

The opposite of \ml{ഇക്കാലത്ത്}, nowadays: \ml{അ}-, that, + \ml{കാലം}, time, in the locative.

\begin{cousinweb}
\textbf{\ml{അക്കാലത്ത്}} (\emph{akkālattŭ}) is built from the far deictic \ml{അ}-, \textquotedblleft{}that\textquotedblright{}, and \ml{കാലം} (\emph{kālaṁ}), \textquotedblleft{}time\textquotedblright{}, from Sanskrit \emph{kāla}, \textquotedblleft{}time\textquotedblright{} (CDIAL 3084), with the locative ending -\ml{ത്ത്}.
\end{cousinweb}

\begin{grammarlens}[title={What you've built}]
One more word for telling what happened.
\end{grammarlens}

\subsection*{Your turn}
\begin{itemize}
\raggedright
  \item \emph{Say it:} \emph{akkālattŭ}
  \item \emph{Say it:} \emph{akkālattŭ}, once more
  \item \emph{Say it:} \emph{akkālattŭ grāmattil basŭ uṇṭāyirunnilla.} — In those days there was no bus in the village.
  \item \emph{Recall:} say the Malayalam for all of a sudden, then the Malayalam for once upon a time, then say \emph{akkālattŭ} again
\end{itemize}

\begin{tcolorbox}[breakable,colback=teal!4,colframe=teal!35!black,title={Before you move on}]
What does \ml{അക്കാലത്ത്} mean? (\textquotedblleft{}In those days\textquotedblright{}.) Say it once more.
\end{tcolorbox}
\section[\texorpdfstring{\ml{അന്നേരം} (annēraṁ)}{അന്നേരം (annēraṁ)}]{\ml{അന്നേരം} (annēraṁ) — at that moment; right then}

\label{lesson:ML-C358-anneram}



\begin{quote}
Before the new one: say the Malayalam for once upon a time, then the Malayalam for in those days.
\end{quote}

\subsection*{You'll want to know: \ml{അന്നേരം}}
\textbf{\ml{അന്നേരം}} — \emph{annēraṁ} — \textquotedblleft{}at that moment; right then\textquotedblright{}.

\begin{quote}
\ml{അന്നേരം} \ml{വാതിലിൽ} \ml{ആരോ} \ml{മുട്ടി}.
\end{quote}

\begin{quote}
\emph{annēraṁ vātilil ārō muṭṭi.}
\end{quote}

\begin{quote}
At that moment someone knocked at the door.
\end{quote}

\ml{അ}-, that, + \ml{നേരം}, time: the exact moment in a story.

\begin{cousinweb}
\textbf{\ml{അന്നേരം}} (\emph{annēraṁ}) is native Dravidian: the far deictic \ml{അ}-, \textquotedblleft{}that\textquotedblright{}, and \ml{നേരം} (\emph{nēraṁ}), \textquotedblleft{}time, season\textquotedblright{} (DEDR 3774).
\end{cousinweb}

\begin{grammarlens}[title={What you've built}]
One more word for telling what happened.
\end{grammarlens}

\subsection*{Your turn}
\begin{itemize}
\raggedright
  \item \emph{Say it:} \emph{annēraṁ}
  \item \emph{Say it:} \emph{annēraṁ}, once more
  \item \emph{Say it:} \emph{annēraṁ vātilil ārō muṭṭi.} — At that moment someone knocked at the door.
  \item \emph{Recall:} say the Malayalam for once upon a time, then the Malayalam for in those days, then say \emph{annēraṁ} again
\end{itemize}

\begin{tcolorbox}[breakable,colback=teal!4,colframe=teal!35!black,title={Before you move on}]
What does \ml{അന്നേരം} mean? (\textquotedblleft{}At that moment; right then\textquotedblright{}.) Say it once more.
\end{tcolorbox}
\section[\texorpdfstring{\ml{തുടർന്ന്} (tuṭarnnŭ)}{തുടർന്ന് (tuṭarnnŭ)}]{\ml{തുടർന്ന്} (tuṭarnnŭ) — following that; subsequently}

\label{lesson:ML-C358-tutarnnu}



\begin{quote}
Before the new one: say the Malayalam for in those days, then the Malayalam for at that moment.
\end{quote}

\subsection*{You'll want to know: \ml{തുടർന്ന്}}
\textbf{\ml{തുടർന്ന്}} — \emph{tuṭarnnŭ} — \textquotedblleft{}following that; subsequently\textquotedblright{}.

\begin{quote}
\ml{ശക്തമായ} \ml{മഴ} \ml{പെയ്തു}, \ml{തുടർന്ന്} \ml{വെള്ളപ്പൊക്കം} \ml{ഉണ്ടായി}.
\end{quote}

\begin{quote}
\emph{śaktamāya maḻa peytu, tuṭarnnŭ veḷḷappokkaṁ uṇṭāyi.}
\end{quote}

\begin{quote}
Heavy rain fell, and following that there was a flood.
\end{quote}

A form of \ml{തുടരുക}, to continue: in what followed. Common in the news.

\begin{cousinweb}
\textbf{\ml{തുടർന്ന്}} (\emph{tuṭarnnŭ}) is native Dravidian, the past participle of \ml{തുടരുക} (\emph{tuṭaruka}), \textquotedblleft{}to continue, to follow on\textquotedblright{} (DEDR 3480). In the news it links one event to the next.
\end{cousinweb}

\begin{grammarlens}[title={What you've built}]
One more word for telling what happened.
\end{grammarlens}

\subsection*{Your turn}
\begin{itemize}
\raggedright
  \item \emph{Say it:} \emph{tuṭarnnŭ}
  \item \emph{Say it:} \emph{tuṭarnnŭ}, once more
  \item \emph{Say it:} \emph{śaktamāya maḻa peytu, tuṭarnnŭ veḷḷappokkaṁ uṇṭāyi.} — Heavy rain fell, and following that there was a flood.
  \item \emph{Recall:} say the Malayalam for in those days, then the Malayalam for at that moment, then say \emph{tuṭarnnŭ} again
\end{itemize}

\begin{tcolorbox}[breakable,colback=teal!4,colframe=teal!35!black,title={Before you move on}]
What does \ml{തുടർന്ന്} mean? (\textquotedblleft{}Following that; subsequently\textquotedblright{}.) Say it once more.
\end{tcolorbox}
\section[\texorpdfstring{\ml{ഉടനീളം} (uṭanīḷaṁ)}{ഉടനീളം (uṭanīḷaṁ)}]{\ml{ഉടനീളം} (uṭanīḷaṁ) — throughout}

\label{lesson:ML-C358-utanilam}



\begin{quote}
Before the new one: say the Malayalam for at that moment, then the Malayalam for following that.
\end{quote}

\subsection*{You'll want to know: \ml{ഉടനീളം}}
\textbf{\ml{ഉടനീളം}} — \emph{uṭanīḷaṁ} — \textquotedblleft{}throughout\textquotedblright{}.

\begin{quote}
\ml{യാത്രയിൽ} \ml{ഉടനീളം} \ml{കുട്ടി} \ml{ഉറങ്ങി}.
\end{quote}

\begin{quote}
\emph{yātrayil uṭanīḷaṁ kuṭṭi uṟaṅṅi.}
\end{quote}

\begin{quote}
The child slept throughout the journey.
\end{quote}

It follows a noun: \ml{വർഷം} \ml{ഉടനീളം}, all through the year.

\begin{cousinweb}
\textbf{\ml{ഉടനീളം}} (\emph{uṭanīḷaṁ}) is native Dravidian. Its second part is \ml{നീളം} (\emph{nīḷaṁ}), \textquotedblleft{}length\textquotedblright{} (DEDR 3692): along the whole length of something. The origin of the first part is not settled.
\end{cousinweb}

\begin{grammarlens}[title={What you've built}]
One more word for telling what happened.
\end{grammarlens}

\subsection*{Your turn}
\begin{itemize}
\raggedright
  \item \emph{Say it:} \emph{uṭanīḷaṁ}
  \item \emph{Say it:} \emph{uṭanīḷaṁ}, once more
  \item \emph{Say it:} \emph{yātrayil uṭanīḷaṁ kuṭṭi uṟaṅṅi.} — The child slept throughout the journey.
  \item \emph{Recall:} say the Malayalam for at that moment, then the Malayalam for following that, then say \emph{uṭanīḷaṁ} again
\end{itemize}

\begin{tcolorbox}[breakable,colback=teal!4,colframe=teal!35!black,title={Before you move on}]
What does \ml{ഉടനീളം} mean? (\textquotedblleft{}Throughout\textquotedblright{}.) Say it once more.
\end{tcolorbox}
